\documentclass[11pt]{article}

\usepackage[a4paper,margin=1in]{geometry}
\usepackage[T1]{fontenc}
\usepackage[utf8]{inputenc}
\usepackage{microtype}
\usepackage{booktabs}
\usepackage{array}
\usepackage{enumitem}
\usepackage{amsmath}
\usepackage{xcolor}
\usepackage{hyperref}
\usepackage{url}

\hypersetup{
  colorlinks=true,
  linkcolor=blue!55!black,
  citecolor=blue!55!black,
  urlcolor=blue!55!black
}
\urlstyle{same}
\emergencystretch=3em
\setlist[itemize]{leftmargin=1.5em}
\setlist[enumerate]{leftmargin=1.5em}

\title{Agent Wiki: Evidence-Closed Local-First Markdown Knowledge Bases for LLM Agents}
\author{Zijun Chu\\
Independent Researcher\\
\texttt{https://github.com/NimaChu/agent-wiki}}
\date{Journal Manuscript Draft, July 2026}

\begin{document}
\maketitle

\begin{abstract}
Large language model (LLM) agents increasingly perform multi-step software and
knowledge-work tasks, but their durable knowledge often remains trapped in chat
transcripts, proprietary note systems, hidden memories, or retrieval indexes
that are difficult for users to inspect. This article presents Agent Wiki, a
local-first, Markdown-native knowledge substrate for LLM coding agents and
human users who need source-grounded, auditable, and low-cost long-term
knowledge management. Agent Wiki separates raw evidence notes from synthesized
wiki pages, treats ordinary wiki links as a computable graph, and introduces an
evidence-closure invariant for processed sources: a source can be considered
processed only when its related wiki targets resolve, the wiki layer links back
to the raw evidence, and explicit follow-up flags are closed. The system also
indexes visual evidence by extracting, mirroring, and inventorying source
images, and it provides an optional read-only dashboard for graph inspection and
maintenance queues. We describe the design, implementation, and reproducibility
workflow, then evaluate the artifact through a single-corpus case study on a
technical documentation vault. On 2026-07-09, the vault contained 1,120 raw
source notes, 113 wiki pages, and 5,353 graph edges; local lint checks reported
zero unresolved links, zero orphaned wiki pages, and zero processed-source
closure violations. Agent Wiki does not replace retrieval-augmented generation;
instead, it provides an inspectable evidence and synthesis layer from which
retrieval, dashboards, and agent workflows can be built.
\end{abstract}

\noindent\textbf{Keywords:}
LLM agents; software engineering; local-first software; Markdown; knowledge
management; evidence-grounded AI; personal knowledge bases; reproducible tools.

\section{Introduction}

LLM agents are moving from short-form assistance toward long-running workflows:
reading documentation, editing source code, maintaining project context,
answering from private corpora, and coordinating multiple local tools. These
workflows require durable knowledge. A useful agent should not repeatedly
rediscover the same documentation, lose the provenance of a claim, or collapse
the distinction between raw evidence and synthesized understanding. Yet many
practical agent memory systems still rely on transient chat history,
application-specific memories, opaque vector indexes, or human-maintained note
vaults that were not designed for automated maintenance.

This creates a gap for independent developers and small teams. They often want
the benefits of an agent-maintained knowledge base without deploying a hosted
database, paying for a specialized knowledge platform, or accepting an opaque
memory layer that cannot be reviewed with ordinary developer tools. They also
need stronger provenance than a flat folder of summaries. When an answer
depends on source material, the system should make it possible to inspect the
original capture, the synthesized wiki page, and the links between them.

Agent Wiki addresses this gap with a deliberately conservative design: local
files first, Markdown first, and evidence first. It is a repository layout, a
set of agent-operating rules, a command line interface, templates, and an
optional dashboard for maintaining a source-grounded local knowledge base. The
raw layer stores source captures and metadata. The wiki layer stores durable
knowledge units, such as concepts, methods, APIs, workflows, entities,
comparisons, and recurring questions. Local scripts scan the Markdown files,
resolve links, build graph statistics, report maintenance queues, and enforce a
processed-source closure rule.

The guiding premise is that agent knowledge should be inspectable before it is
automated. Agent Wiki therefore treats the graph as a derived artifact rather
than the source of truth. A user can inspect, diff, search, copy, version,
publish, or delete the underlying files without needing the dashboard or a
database server. This local-first posture follows the principle that users
should retain ownership of their data even when cloud services or agents are
used around it \cite{kleppmann2019localfirst}.

This article makes four contributions:

\begin{enumerate}
  \item It presents a file-native architecture for agent-maintained knowledge
  bases that separates raw evidence from synthesized wiki knowledge.
  \item It formalizes an evidence-closure invariant for processed source notes
  and implements local checks for unresolved links, orphaned wiki pages,
  invalid relations, and processed-gate failures.
  \item It introduces an image-evidence workflow that mirrors and inventories
  visual source material so screenshots, diagrams, and UI states can remain
  available as local evidence.
  \item It reports a case-study evaluation over a technical documentation
  vault, showing that local health checks can preserve graph consistency across
  more than one thousand raw source notes.
\end{enumerate}

\section{Background and Motivation}

\subsection{Agent Workflows Need External Memory}

LLMs are increasingly used as controllers that interleave reasoning, tool use,
and environment feedback. ReAct showed that language models can combine
reasoning traces with actions to solve interactive tasks \cite{yao2023react}.
Software-engineering agents extend this idea to repositories, issue trackers,
tests, terminals, and code search. Benchmarks and systems such as SWE-bench and
SWE-agent demonstrate that code-oriented agents can attempt real software tasks
when equipped with appropriate interfaces \cite{jimenez2024swebench,
yang2024sweagent}.

As these workflows lengthen, memory becomes a practical systems problem. A
coding agent may need to remember local conventions, documentation facts,
decisions, evidence, screenshots, and unresolved follow-ups. Keeping this
knowledge only in chat history is fragile. Storing it only in an embedding
index improves retrieval but can make provenance and maintenance harder to
inspect. Storing it only in a human note-taking application may help the human
but leaves the agent without a disciplined operating model.

\subsection{RAG Is Complementary but Not Sufficient}

Retrieval-augmented generation (RAG) connects a language model to external
documents at answer time and has become a standard pattern for knowledge-
intensive generation \cite{lewis2020rag}. However, RAG does not by itself
specify how source captures should be preserved, how synthesized concepts
should evolve, how evidence links should be audited, or how an agent should
decide that a source has been fully processed. RAG optimizes retrieval; Agent
Wiki optimizes the preparation and maintenance of the corpus that retrieval may
later consume.

This distinction is important for small local projects. A user may want a
knowledge base that can be read without a model, maintained with ordinary text
tools, and later embedded or indexed if needed. Agent Wiki is therefore not an
alternative to RAG in the sense of answer-time retrieval. It is a staging and
audit layer for source-grounded knowledge.

\subsection{Personal Knowledge Tools Are Not Agent Protocols}

Markdown note systems such as Obsidian and Logseq show the value of local
files, backlinks, and graph-oriented thinking \cite{obsidian,logseq}. Evergreen
note practices emphasize durable, atomic notes that become more useful through
linking and revision \cite{matuschakEvergreen}. Agent Wiki adopts these habits
but narrows them for agentic operation. It defines a repository structure,
templates, status fields, closure rules, and scripts that an LLM coding agent
can execute consistently.

The difference is procedural rather than cosmetic. Agent Wiki is not merely a
folder of Markdown files. It specifies what counts as raw evidence, what counts
as synthesized wiki knowledge, when a source may be marked processed, which
links must resolve, which relation hints are supported, and which dashboard
artifacts are read-only.

\section{Design Goals}

Agent Wiki is shaped by six design goals.

\paragraph{Local ownership.}
The user should own source captures, synthesized pages, image inventories,
snapshots, and maintenance state as local files. Public project code and private
knowledge can be separated by Git policy.

\paragraph{Agent-native operation.}
The repository should be legible to coding agents. The project includes
workspace rules, templates, and root npm scripts so agents can capture sources,
search, maintain, lint, repair links, and answer from the knowledge base without
depending on globally installed skills.

\paragraph{Evidence before synthesis.}
Source captures should remain available after synthesis. The wiki layer is a
compiled knowledge layer, not a replacement for evidence.

\paragraph{Low operational cost.}
The system should not require a hosted database, a paid subscription, an
embedding service, an always-on dashboard, or a separate editor runtime.

\paragraph{Auditable maintenance.}
Maintenance state should be computable from files. Broken links, orphaned wiki
pages, unresolved related targets, stale sources, follow-up flags, and
processed-gate failures should be visible to local commands.

\paragraph{Incremental extensibility.}
The file-native substrate should not prevent future indexing, dashboards,
retrieval tools, or structured agent APIs. Derived artifacts should be
rebuildable from Markdown.

\section{System Overview}

Agent Wiki is organized as a self-contained repository. Table~\ref{tab:layout}
summarizes the major directories.

\begin{table}[t]
\centering
\begin{tabular}{>{\raggedright\arraybackslash}p{0.28\linewidth}p{0.62\linewidth}}
\toprule
Directory & Role \\
\midrule
\path|raw/| & Source notes, snapshots, evidence metadata, and image inventories. \\
\path|wiki/| & Durable synthesized knowledge pages. \\
\path|templates/| & Reusable raw-source and wiki-page templates. \\
\path|src/| & Canonical command line source for capture, search, lint, repair, and graph generation. \\
\path|scripts/| & Backwards-compatible command launchers and optional local helpers. \\
\path|tools/wiki-dashboard/| & Optional read-only graph dashboard. \\
\bottomrule
\end{tabular}
\caption{Agent Wiki repository layout.}
\label{tab:layout}
\end{table}

\subsection{Raw Evidence Layer}

The raw layer is the factual source layer. A raw note stores bibliographic and
capture metadata such as title, author, date, URL, captured date, type, status,
snapshot path, content hash, capture method, source quality, tags, related wiki
targets, and follow-up state. The body can preserve captured text, extracted
claims, image inventories, candidate wiki links, and processing notes.

Raw notes are intentionally not treated as disposable summaries. A single raw
source may update many wiki pages, and a single wiki page may link back to many
raw sources. This design prevents a source from being collapsed into one
summary page when it contains multiple reusable ideas.

\subsection{Wiki Layer}

The wiki layer is the compiled knowledge layer. It favors one durable knowledge
unit per page: a concept, entity, API, method, workflow, comparison, or
recurring question. Pages use Obsidian-compatible links such as
\texttt{[[Page Name]]} and \texttt{[[Page Name|label]]}. Navigation files such
as \path|wiki/index.md| and \path|wiki/log.md| are treated as maintenance
records rather than graph knowledge nodes.

Wiki pages may include typed relation hints from a deliberately small
vocabulary: \texttt{supports}, \texttt{challenges}, \texttt{related\_to},
\texttt{applies\_to}, \texttt{company\_of}, and \texttt{product\_of}. The small
set keeps relations easy to validate and avoids turning the vault into an
uncontrolled ontology project.

\subsection{Command Line Interface}

The root command line interface exposes local workflows through npm scripts. The
main entry point delegates to specialized scripts for capture, status, lint,
garden review, link repair, search, image indexing, IMA import, and dashboard
generation. Table~\ref{tab:commands} lists the primary commands.

\begin{table}[t]
\centering
\begin{tabular}{>{\ttfamily}p{0.38\linewidth}p{0.52\linewidth}}
\toprule
\normalfont Command & Purpose \\
\midrule
npm run wiki:status & Report vault graph and queue statistics. \\
npm run wiki:lint & Validate unresolved links, metadata, and processed-source closure. \\
npm run wiki:garden & Show maintenance queues and weakly connected pages. \\
npm run wiki:repair-links & Assist link repair. \\
npm run wiki:search -- "query" & Search titles, tags, aliases, and excerpts. \\
npm run wiki:capture & Create raw source notes. \\
npm run wiki:images & Extract, mirror, and index image evidence. \\
npm run dashboard:open & Refresh graph, run the local dashboard, and open the browser. \\
\bottomrule
\end{tabular}
\caption{Primary local workflows exposed as root npm scripts.}
\label{tab:commands}
\end{table}

\section{Evidence-Closure Model}

The key invariant in Agent Wiki is evidence closure for processed raw sources.
Let $R$ be the set of raw source notes and $W$ the set of wiki pages. For a raw
source $r \in R$, let $\mathrm{related}(r)$ be the list of wiki targets declared
in the source metadata, and let $\mathrm{resolve}(x)$ map a wiki-link target to
an existing node or to failure. Let $\mathrm{links}(w)$ be the set of links
from a wiki page $w$.

A processed source is considered closed only if:

\[
\begin{aligned}
\mathrm{processed}(r) \Rightarrow\;&
|\mathrm{related}(r)| > 0\\
&\land\ \forall t \in \mathrm{related}(r),\ \mathrm{resolve}(t) \neq \bot\\
&\land\ \exists w \in W,\ w \in \mathrm{resolve}(\mathrm{related}(r))
\land r \in \mathrm{links}(w)\\
&\land\ \neg \mathrm{needsFollowup}(r).
\end{aligned}
\]

This invariant is intentionally modest. It does not prove that every claim in a
wiki page is correct. It does not prove that the source has been exhaustively
processed. It does, however, prevent a common agent-maintenance failure: a
source is marked done but is unreachable from the durable wiki layer, or the
wiki layer has no backlink to the evidence that supposedly supports it.

The lint workflow reports processed-source issues using four reason codes:
\texttt{missing-related}, \texttt{unresolved-related},
\texttt{missing-wiki-backlink}, and \texttt{explicit-followup}. This turns a
subjective maintenance label into a reviewable local property.

\section{Image Evidence}

Many sources contain visual evidence: UI states, screenshots, diagrams,
configuration panels, charts, and workflows. Ordinary text-only capture can
lose this evidence or leave it dependent on remote image URLs. Agent Wiki
therefore includes a dedicated image workflow.

The image script scans a raw note and selected snapshots for Markdown image
syntax and HTML \texttt{img} tags. It resolves relative URLs against the source
URL, deduplicates images, mirrors remote images into
\path|raw/assets/<source-note>/|, records local paths, content type, byte
counts, source context, and download errors, writes a machine-readable
\path|image-index.json|, and updates the raw note's \texttt{Images} section
with a compact inventory table.

Agents are instructed to promote only useful visual evidence into wiki pages.
The exhaustive inventory remains in the raw layer. This keeps answers concise
while preserving the complete visual record when a user needs to audit a
screenshot or diagram.

\section{Graph Generation and Dashboard}

The scanner walks Markdown files under \path|raw/|, \path|wiki/|,
\path|templates/|, and archives. It parses simple frontmatter, extracts body
links and frontmatter links, resolves targets by id, title, basename, and
alias, and emits nodes, edges, unresolved links, typed relations, invalid
relations, incoming counts, outgoing counts, and node excerpts.

The dashboard is a read-only visualization surface. It displays global and
local graph views, a wiki-only knowledge view by default, evidence drill-down
for one wiki page and its directly linked raw sources, search, group and status
filters, selected-node links and backlinks, tags, raw/wiki/link counts, and
health indicators for inbox, follow-up, stale, broken-link, and processed-gate
queues.

The dashboard intentionally does not write to the vault. Write operations remain
explicit Markdown edits or command line workflows, which makes the graph a
derived artifact rather than a hidden source of truth.

\section{Implementation}

Agent Wiki is implemented with Node.js scripts and a Vite/React dashboard. The
public software artifact is small and inspectable: the root package exposes the
core scripts, and the dashboard is an optional local frontend. The system has no
required database server, no required cloud backend, and no required Obsidian
runtime.

The main scanner is implemented in \path|src/wiki-lib.mjs|. It provides
frontmatter parsing, wiki-link extraction, link resolution, relation parsing,
vault scanning, graph statistics, and processed-source issue detection. The
command dispatcher in \path|src/karpathy-wiki.mjs| maps user-facing npm
scripts to implementation modules. Image indexing is implemented in
\path|src/image-assets.mjs|. Dashboard generation is implemented under
\path|tools/wiki-dashboard/|.

The implementation chooses simple structured parsing over ad hoc shell
pipelines. This matters because the same rules are used by human users, coding
agents, lint commands, and dashboard generation. A link that resolves in the
dashboard should resolve in lint and search as well.

\section{Evaluation Method}

The evaluation is a case-study evaluation of the artifact as a local
maintenance substrate. The goal is not to show that Agent Wiki improves model
answer quality over all RAG systems. Instead, the evaluation asks whether the
system can maintain inspectable graph health over a documentation-scale vault
using local commands.

We ask three research questions:

\begin{description}
  \item[RQ1:] Can the vault maintain link and evidence-closure health over more
  than one thousand source notes?
  \item[RQ2:] Can health be checked by local, reproducible commands rather than
  by manual dashboard inspection?
  \item[RQ3:] What practical limitations remain for journal-level claims and
  future evaluation?
\end{description}

The case-study corpus is a local technical documentation vault centered on
Autodesk FlexSim documentation. The vault was maintained using Agent Wiki
workflows and checked on 2026-07-09 with:

\begin{verbatim}
npm run wiki:status
npm run wiki:lint
\end{verbatim}

The commands report graph statistics, unresolved links, invalid relations,
processed-source issues, orphaned wiki pages, weak wiki pages, missing claim
sources, and missing metadata fields.

\section{Case-Study Results}

Table~\ref{tab:case} reports the vault state on 2026-07-09. The local vault
contained 1,244 graph nodes, 5,353 graph edges, 1,120 raw source notes, and 113
wiki pages. Lint reported no unresolved links, invalid relation hints,
processed-source issues, orphaned wiki pages, weak wiki pages, missing claim
sources, missing frontmatter, missing status fields, or missing type fields.

\begin{table}[t]
\centering
\begin{tabular}{lr}
\toprule
Metric & Value \\
\midrule
Total nodes & 1,244 \\
Graph edges & 5,353 \\
Raw source notes & 1,120 \\
Wiki pages & 113 \\
Pending raw sources & 0 \\
Inbox raw sources & 0 \\
IMA pointer notes & 0 \\
Processed raw sources & 1,119 \\
Needs follow-up & 0 \\
Stale sources & 0 \\
Unresolved links & 0 \\
Invalid relation hints & 0 \\
Orphaned wiki pages & 0 \\
Processed-source closure issues & 0 \\
\bottomrule
\end{tabular}
\caption{Case-study vault health reported by Agent Wiki local commands on
2026-07-09.}
\label{tab:case}
\end{table}

These results answer RQ1 and RQ2 for the case-study setting. The vault remained
mechanically healthy under the implemented checks despite containing more than
one thousand raw source notes. The absence of processed-source closure issues is
particularly important: it indicates that processed source notes had resolved
related wiki targets and corresponding wiki backlinks, rather than being
isolated summaries.

The result should be interpreted as a systems-health result, not an answer-
quality benchmark. Link health and closure health are necessary but not
sufficient for factual correctness. They provide an auditable substrate on which
answer quality can later be evaluated.

\section{Discussion}

\subsection{Why Markdown First?}

Markdown gives Agent Wiki a low-friction shared representation for humans,
agents, Git, search tools, and dashboards. It can be edited in a terminal or a
graphical editor, diffed in Git, indexed by scripts, and copied without export
steps. This does not make Markdown semantically rich by itself. Agent Wiki adds
structure through conventions: frontmatter fields, source statuses, supported
relation hints, raw/wiki separation, and closure checks.

\subsection{Why Not Start with Embeddings?}

Embeddings are useful for retrieval, but they are a derived index. Starting with
an embedding database can hide source quality, duplicate captures, weak
synthesis, broken backlinks, and unresolved follow-ups. Agent Wiki starts from
inspectable source and synthesis files. An embedding index can be added later
from either layer.

\subsection{What Agents Gain}

An agent using Agent Wiki gains operational handles rather than only documents.
It can search the vault, inspect the index, read source-backed wiki pages,
detect broken links, process inbox items, repair evidence closure, and include
local visual evidence in answers. The repository itself tells the agent how to
behave, which reduces reliance on hidden application memory.

\subsection{What Users Gain}

Users gain a private knowledge base whose maintenance state can be audited
without trusting the agent. If an answer cites a wiki page, the user can inspect
that page. If the page cites raw evidence, the user can inspect the source note,
snapshot, or image inventory. If the vault claims a source is processed, the
user can run the same lint checks.

\section{Threats to Validity}

\paragraph{Single-corpus evaluation.}
The case study uses one local technical documentation corpus. The result may not
generalize to highly dynamic web sources, collaborative writing, legal
materials, scientific PDFs, or multimodal corpora with large binary assets.

\paragraph{Single-developer workflow.}
The current artifact was developed and maintained primarily by one independent
developer with agent assistance. Multi-user editing, merge conflicts, role-based
review, and permission models require further study.

\paragraph{Health is not correctness.}
The reported metrics measure graph consistency and closure invariants. They do
not prove that every wiki claim faithfully summarizes its sources or that agent
answers are factually superior to RAG baselines.

\paragraph{Public artifact versus local corpus.}
The software artifact is public, but private or bulky raw knowledge can remain
local by design. This improves privacy but means that some case-study vault
content may require regeneration from public sources or a separately published
snapshot for independent replication.

\paragraph{Simple parsing model.}
The current frontmatter parser and link resolver intentionally cover a simple
Markdown subset. This improves portability but may not support every convention
used by larger static-site generators or note-taking systems.

\section{Future Work}

Future work falls into five areas. First, Agent Wiki could expose a Model
Context Protocol server so agents can query vault state through structured APIs
while preserving Markdown as the source of truth. Second, optional embedding
indexes could be generated from raw notes, wiki pages, or both, enabling
comparative evaluation of graph-guided search and vector retrieval. Third,
evidence closure could move from page-level backlinks to claim-level citation
checks. Fourth, image evidence could be enriched with OCR, diagram captions, and
visual similarity search. Fifth, future studies should compare answer quality,
maintenance effort, and user trust against chat-only memory, note-only
workflows, and RAG-first baselines.

\section{Conclusion}

Agent Wiki shows that an LLM-agent memory substrate can begin as a local
Markdown repository rather than a hosted service or hidden vector database. By
separating raw evidence from synthesized wiki knowledge, enforcing processed-
source closure, indexing images as evidence, and deriving graph health from
plain files, the system offers a practical bridge between personal knowledge
management, source-grounded agent workflows, and future retrieval pipelines.
The case-study vault demonstrates that the approach can maintain local graph
health across more than one thousand source notes. The broader claim is modest
but important: before agent knowledge becomes more automated, it should become
more inspectable.

\section*{Availability}

The Agent Wiki software artifact is available at
\url{https://github.com/NimaChu/agent-wiki}. The archived software release used
for this manuscript is available through Zenodo at
\url{https://doi.org/10.5281/zenodo.21252445}. The manuscript draft was prepared
against repository release \texttt{v0.2.1} and local vault status checked on
2026-07-15. The public release contains the reusable software, paper materials,
and a metadata-only FlexSim 2026 case-study subset; private or bulky raw
knowledge vaults can be kept local by design.

\section*{CRediT Author Statement}

Zijun Chu: Conceptualization, methodology, software, validation, investigation,
writing--original draft, writing--review and editing, visualization, project
administration.

\section*{Declaration of Competing Interest}

The author declares no known competing financial interests or personal
relationships that could have appeared to influence the work reported in this
paper.

\section*{Funding}

This research received no specific grant from any funding agency in the public,
commercial, or not-for-profit sectors.

\section*{Acknowledgements}

The author thanks the open-source communities whose tools and documentation
make local-first software development practical.

\begin{thebibliography}{99}

\bibitem{kleppmann2019localfirst}
Martin Kleppmann, Adam Wiggins, Peter van Hardenberg, and Mark McGranaghan.
\newblock Local-first software: You own your data, in spite of the cloud.
\newblock In \emph{Proceedings of the 2019 ACM SIGPLAN International Symposium
on New Ideas, New Paradigms, and Reflections on Programming and Software},
pages 154--178, 2019.

\bibitem{yao2023react}
Shunyu Yao, Jeffrey Zhao, Dian Yu, Nan Du, Izhak Shafran, Karthik
Narasimhan, and Yuan Cao.
\newblock ReAct: Synergizing reasoning and acting in language models.
\newblock In \emph{International Conference on Learning Representations}, 2023.
\newblock \url{https://arxiv.org/abs/2210.03629}.

\bibitem{jimenez2024swebench}
Carlos E. Jimenez, John Yang, Alexander Wettig, Shunyu Yao, Kexin Pei,
Ofir Press, and Karthik Narasimhan.
\newblock SWE-bench: Can language models resolve real-world GitHub issues?
\newblock In \emph{International Conference on Learning Representations}, 2024.
\newblock \url{https://arxiv.org/abs/2310.06770}.

\bibitem{yang2024sweagent}
John Yang, Carlos E. Jimenez, Alexander Wettig, Kilian Lieret, Shunyu Yao,
Karthik Narasimhan, and Ofir Press.
\newblock SWE-agent: Agent-computer interfaces enable automated software
engineering.
\newblock arXiv preprint, 2024.
\newblock \url{https://arxiv.org/abs/2405.15793}.

\bibitem{lewis2020rag}
Patrick Lewis, Ethan Perez, Aleksandra Piktus, Fabio Petroni, Vladimir
Karpukhin, Naman Goyal, Heinrich Kuttler, Mike Lewis, Wen-tau Yih, Tim
Rocktaschel, Sebastian Riedel, and Douwe Kiela.
\newblock Retrieval-augmented generation for knowledge-intensive NLP tasks.
\newblock In \emph{Advances in Neural Information Processing Systems}, 2020.
\newblock \url{https://arxiv.org/abs/2005.11401}.

\bibitem{packer2023memgpt}
Charles Packer, Sarah Wooders, Kevin Lin, Vivian Fang, Shishir G. Patil,
Ion Stoica, and Joseph E. Gonzalez.
\newblock MemGPT: Towards LLMs as operating systems.
\newblock arXiv preprint, 2023.
\newblock \url{https://arxiv.org/abs/2310.08560}.

\bibitem{matuschakEvergreen}
Andy Matuschak.
\newblock Evergreen notes.
\newblock \url{https://notes.andymatuschak.org/}.

\bibitem{obsidian}
Obsidian.
\newblock Obsidian.
\newblock \url{https://obsidian.md/}.

\bibitem{logseq}
Logseq.
\newblock Logseq.
\newblock \url{https://logseq.com/}.

\bibitem{khoj}
Khoj AI.
\newblock Khoj documentation.
\newblock \url{https://docs.khoj.dev/}.

\end{thebibliography}

\end{document}
